\chapter{Hecke 特征性质}\label{ch:2}
% 原PDF第11页下半部
几何 Langlands 理论的自守一侧围绕 $\Bun_G$ 展开；这是 $X$ 上 $G$ 丛的模叠。有了 $\Bun_G$，就可以借助在曲线上某一点修改 $G$ 丛的对应来表述 Hecke 特征性质。

以下首先回顾定义 $\Bun_G$ 所用的预叠与叠的语言，记录其基本有限性性质，并将它的 $k$-点与通常的阿代尔双商联系起来。
\section{预叠与叠}\label{sec:2.1}
\srcpage{12}\index{预叠|hyperpage}\index{叠|hyperpage}\index{下降|hyperpage}
\begin{definition}{}{2.1.1}
$k$ 上的预叠是一个伪函子
\[y:\operatorname{Aff}_k^{\mathrm{op}}\longrightarrow\operatorname{Grpd}.\]
因此，每个仿射 $k$-概形 $S$ 给出一个群胚 $y(S)$，每个映射 $f:S_1\to S_2$ 给出拉回函子 $f^*:y(S_2)\to y(S_1)$，满足通常的相容性条件。若目标范畴是 $\operatorname{Set}$，这就是通常的预层概念。
\end{definition}

设 $\tau$ 是 $\operatorname{Aff}_k$ 上的 Grothendieck 拓扑，例如 Zariski、étale、光滑、fppf 或 fpqc 拓扑。
\begin{definition}{}{2.1.2}
若预叠 $y$ 对 $\tau$-覆盖满足下降，则称它为 $\tau$-叠。等价地，对每个 $\tau$-覆盖 $S^0\to S$ 及其 Čech 神经 $S^\bullet$，自然函子
\[y(S)\longrightarrow\lim y(S^\bullet)\]
是群胚的等价。
\end{definition}
\begin{note}\numberednote{2.1.3}\textbf{（分类叠）}
对群概形 $G/k$，$\tau$-分类叠 $B_\tau G$ 是商预叠 $\mathrm{pt}/G$ 的叠化。具体地，
\[B_\tau G(S)=\{S\text{ 上 }\tau\text{-局部平凡的 }G\text{ 挠子}\}.\]
它可以写为 $\operatorname{Stk}_k$ 中条形构造的叠论余极限：
\[B_\tau G\simeq\operatorname*{colim}(\mathrm{pt}\leftarrow G\leftleftarrows G\times G\cdots).\]
若 $G$ 在 $k$ 上光滑，则挠子在 étale 拓扑下就已局部平凡，因此
\[B_{\mathrm{et}}G=B_{\mathrm{sm}}G=B_{\mathrm{fppf}}G=B_{\mathrm{fpqc}}G.\]
\end{note}
\section{\texorpdfstring{$G$}{G} 丛的模叠}\label{sec:2.2}
\index{模叠!G丛|hyperpage}\index{Artin 叠|hyperpage}
设 $G/k$ 是光滑仿射代数群（在下文几何 Langlands 的应用中为连通约化群），$p:X\to S$ 是射影、平坦且有限表现的态射。定义
\[\Bun_{G,X/S}:=\operatorname{Maps}_S(X,B_{\mathrm{et}}G).\]
对任意 $S'\to S$，
\[\Bun_{G,X/S}(S')=\Hom(X\times_S S',B_{\mathrm{et}}G)=\{X\times_S S'\text{ 上的 étale }G\text{ 挠子}\}.\]
当 $S=\Spec k$ 时，简记为 $\Bun_G$。
\begin{theorem}{}{2.2.1}
叠 $\Bun_{G,X/S}$ 是在 $S$ 上局部有限表现的 Artin 叠，且具有仿射对角。更明确地，
\[\Delta_{\Bun_{G,X/S}}:\Bun_{G,X/S}\longrightarrow\Bun_{G,X/S}\times_S\Bun_{G,X/S}\]
是仿射态射，且 $\Bun_{G,X/S}$ 有开子叠穷竭 $\Bun_{G,X/S}=\bigcup_r U_r$，使每个 $U_r\to\Bun_{G,X/S}$ 都是开嵌入，每个 $U_r$ 都有光滑图册 $V_r\to U_r$，且 $V_r\to S$ 有限表现。
\end{theorem}
\begin{note}\numberednote{2.2.2}
若 $G/k$ 光滑且 $X\to S$ 是相对曲线，则 $\Bun_{G,X/S}\to S$ 光滑。在上述穷竭中，可以选择使 $V_r\to S$ 光滑的图册。
\end{note}
\srcpage{13}
\begin{note}\numberednote{2.2.3}\textbf{（$G=\GL_n$）}
固定 $X$ 上的丰沛线丛 $\cO(1)$。对 $r\gg0$，令 $U_r\subset\Bun_n$ 为由满足
\[\mathrm R^ip_*(\cF(r))=0\quad(i>0),\qquad p^*p_*(\cF(r))\twoheadrightarrow\cF(r)\]
的秩 $n$ 向量丛 $\cF$ 组成的开子叠。开子叠 $U_r$ 覆盖 $\Bun_n$。不过，$\Bun_G$ 通常不拟紧；对 $G=\GL_n$，它不连通，甚至一个固定的分支 $\Bun_G^\lambda$ 也未必拟紧。
\end{note}

以下阿代尔描述是 $\Bun_G$ 在函数论中的体现。令 $K=k(X)$，$K_x$ 为 $x\in X$ 处的完备局部域，并置 $\A_X=\prod_x'K_x$，$\cO_X=\prod_x\cO_x$。
\begin{theorem}{$k$-点上的阿代尔一致化}{2.2.4}
假设 $\Spec K$ 上、每个 $\Spec K_x$ 上以及每个有限扩张 $k'/k$ 的 $\Spec k'$ 上的所有 $G$ 挠子都是平凡的。则采用右挠子的约定，自然的粘合构造将 $\Bun_G$ 的 $k$-点同构类等同于双商
\[G(\cO_X)\backslash G(\A_X)/G(K).\]
采用左挠子时，同一陈述写成通常的相反次序 $G(K)\backslash G(\A_X)/G(\cO_X)$。
\end{theorem}
\begin{note}\numberednote{2.2.5}
当 $k$ 有限时，平凡性假设特别适用于单连通半单群，例如 $\SL_2$，以及 $\GL_n$。当 $k=\bar k$ 且特征为 $0$ 时，它适用于连通线性代数群 $G$；当 $k=\bar k$ 且特征为 $p$ 时，它适用于连通约化群 $G$。
\end{note}

本章其余部分介绍整体与局部 Hecke 叠，回顾定义 Hecke 函子的 Satake 层，然后专门讨论 $\GL_n$，以获得微小 Hecke 特征性质。
\section{整体 Hecke 叠}\label{sec:2.3}\index{Hecke 叠!整体|hyperpage}
从现在起假设 $X$ 是 $k$ 上的光滑曲线。整体 Hecke 叠记录 $G$ 丛在 $X$ 上一个移动点处的一次修改。
\begin{definition}{}{2.3.1}
叠 $\Hecke_{G,X}$ 定义为
\[
\begin{aligned}
\Hecke_{G,X}(S)=\{(x,P,P',\beta)\mid{}&x:S\to X,\ P,P'\in\Bun_G(S),\\
&\beta:P|_{X\times S-\Gamma_x}\xrightarrow{\sim}P'|_{X\times S-\Gamma_x}\},
\end{aligned}
\]
其中 $\Gamma_x\subset X\times S$ 是 $x$ 的图像。它有如下对应图：
\origdiagram{p013-diagram.pdf}{0.5146}
这里，$\overleftarrow h$ 记住 $P$，$\overrightarrow h$ 记住 $P'$，$\supp$ 记住点 $x$。
\end{definition}
\section{仿射 Grassmann 簇与局部模型}\label{sec:2.4}\index{仿射 Grassmann 簇|hyperpage}
$X$ 上的 Beilinson–Drinfeld 仿射 Grassmann 簇是 Hecke 叠的局部模型：
\[
\begin{aligned}
\Gr_{G,X}(S)=\{(x,P,\beta)\mid{}&x:S\to X,\ P\text{ 为 }X\times S\text{ 上的 }G\text{ 挠子},\\
&\beta:P|_{X\times S-\Gamma_x}\xrightarrow{\sim}P^0|_{X\times S-\Gamma_x}\},
\end{aligned}
\]
其中 $P^0$ 是平凡 $G$ 挠子。
\begin{theorem}{}{2.4.1}
$\Gr_{G,X}$ 是 ind-局部有限型的 ind-概形。若 $G$ 约化，则 $\Gr_{G,X}$ 在 $X$ 上 ind-射影。
\end{theorem}
\srcpage{14}

固定 $x:S\to X$。令 $D_x=\widehat{(X\times S)}_{\Gamma_x}$ 为沿图像的形式圆盘，$D_x^\circ=D_x\setminus\Gamma_x$ 为穿孔形式圆盘。定义正环路预叠和全环路预叠：
\[L^+G_X(S)=\{(x,\gamma)\mid\gamma:D_x\to G\},\qquad LG_X(S)=\{(x,\gamma)\mid\gamma:D_x^\circ\to G\}.\]
若 $S=\Spec k$ 且 $x$ 有参数 $t$，则 $D_x=\Spec k[[t]]$，$D_x^\circ=\Spec k((t))$。分类叠 $BL^+G_X$ 和 $BLG_X$ 分别分类 $D_x$ 与 $D_x^\circ$ 上的 $G$ 挠子。仿射 Grassmann 簇是商
\[\Gr_{G,X}\simeq LG_X/L^+G_X.\]

为了从这个局部模型过渡到整体修改，引入无限级丛
\[\Bun_G^\infty(S)=\{(x,P,\alpha)\mid x:S\to X,\ P\in\Bun_G(S),\ \alpha:P|_{D_x}\xrightarrow{\sim}P^0|_{D_x}\}.\]
它是 $X\times\Bun_G$ 上的 $L^+G_X$ 挠子。因此，从直观上以及精确的商意义上都有
\[\Hecke_{G,X}\simeq\Bun_G^\infty\times^{L^+G_X}\Gr_{G,X}\simeq\Gr_{G,X}\widetilde\times\Bun_G.\]
从而 $\Hecke_{G,X}$ 是 ind-局部有限型的 ind-Artin 叠。

\begin{transnote}{含糊点拨}
「无限级」与维数无关：$\alpha$ 定义在形式圆盘 $D_x=\Spec k[[t]]$（而非某阶 $D_x^{(n)}=\Spec k[t]/(t^n)$）上，因此 $\Bun_G^\infty$ 是有限级挠子 $\Bun_G^{(n)}$ 的逆向极限，只是 ind-对象，并不对应概形。扭曲积 $\Bun_G^\infty\times^{L^+G_X}\Gr_{G,X}$ 的含义是：一点由「$L^+G_X$-挠子 $\alpha$」与「$\Gr$ 的一点，即 $D_x^\circ$ 上的 $G$-丛配平凡化」组成，模去 $L^+G_X$ 的右作用后，正是「修改前的 $P$、修改后的 $P'$、以及 $D_x^\circ$ 上的同构」这三元组，也就是 $\Hecke_{G,X}$ 的定义。所以第二个同构并非巧合：$\Gr_{G,X}\widetilde\times\Bun_G$ 与扭曲积是同一叠的两种写法。
[来源：nLab "affine Grassmannian"，https://ncatlab.org/nlab/show/affine+Grassmannian（Last revised 2023-08-15；2026-11-17 访问）；并见原书 p.14 对 $D_x$ 与 $L^+G_X$ 的约定。]
\end{transnote}
\section{局部 Hecke 叠与几何 Satake}\label{sec:2.5}\index{Hecke 叠!局部|hyperpage}\index{几何 Satake|hyperpage}
对固定点 $x\in X$，记 $L^+G_x=G(\cO_x)$，$LG_x=G(K_x)$。局部 Hecke 叠为
\[\Hecke_{G,x}^{\mathrm{loc}}:=[L^+G_x\backslash LG_x/L^+G_x]\simeq[L^+G_x\backslash\Gr_{G,x}].\]
将一个修改在 $x$ 处完备化，便得到从整体 Hecke 叠到局部 Hecke 叠的相对位置映射
\[m_x:\Hecke_{G,x}\longrightarrow\Hecke_{G,x}^{\mathrm{loc}}.\]

几何 Satake 将局部 Hecke 叠上的球面层范畴等同于 Langlands 对偶群的表示范畴：
\[\operatorname{Perv}_{L^+G_x}(\Gr_{G,x})\simeq\operatorname{Perv}(\Hecke_{G,x}^{\mathrm{loc}})\simeq\Rep(\checkG).\]
对 $V\in\Rep(\checkG)$，以 $\operatorname{Sat}_V$ 表示 $\Hecke_{G,x}^{\mathrm{loc}}$ 上相应的反常层；在 $X$ 上的整体族中，也用同一记号表示相应的相对 Satake 层。

\begin{transnote}{背景直觉}
这句话的内容比「认同」更强：$\operatorname{Perv}_{L^+G}(\Gr_G)$ 上由卷积（即融合积）给出的张量结构，与 $\Rep(\checkG)$ 的张量结构等价；对应把 $\operatorname{IC}(\overline{\Gr^\lambda})$ 送到最高权表示 $V^\lambda$，其中 $\lambda$ 跑优势余权、$\overline{\Gr^\lambda}$ 是 $\Gr^\lambda$ 的闭包。$\Hecke^{\mathrm{loc}}$ 上的球面层与 $\Gr$ 上的 $L^+G$-等变反常层是同一批对象，「球面」指按 $L^+G$ 的右作用取不变。$\ell$-进情形还须用 Weil 结构处理 Frobenius 归一化，Betti 与 de Rham 情形则直接是反常层。这句话同时给出了 $\operatorname{Sat}_V$ 的定义：$V$ 的不可约分解对应相应轨道 IC 层的直和。
[来源：nLab "geometric Satake equivalence"，https://ncatlab.org/nlab/show/geometric+Satake+equivalence（Last revised 2024-07-22；2026-11-17 访问）；并见 [FGKV98]（JAMS 11 (1998), 451–484）。]
\end{transnote}
\section{Hecke 函子与特征层}\label{sec:2.6}\index{Hecke 函子|hyperpage}
对 $V\in\Rep(\checkG)$，整体 Hecke 函子为
\[
\begin{aligned}
H_V:D(\Bun_G)&\longrightarrow D(X\times\Bun_G),\\
\cF&\longmapsto(\supp,\overrightarrow h)_!(m^*\operatorname{Sat}_V\otimes\overleftarrow h^*\cF).
\end{aligned}
\]
在固定点 $x\in X$ 处，它限制为
\[
\begin{aligned}
H_{V,x}:D(\Bun_G)&\longrightarrow D(\Bun_G),\\
\cF&\longmapsto\overrightarrow h_{x,!}(m_x^*\operatorname{Sat}_V\otimes\overleftarrow h_x^*\cF).
\end{aligned}
\]
$\operatorname{Sat}_V$ 的反常归一化包括通常的移位与扭曲。

设 $E_{\checkG}$ 是 $X$ 上的 $\checkG$-局部系统。对 $V\in\Rep(\checkG)$，令 $V_E$ 为 $X$ 上相应的局部系统。
\begin{definition}{}{2.6.1}
若对每个 $V\in\Rep(\checkG)$，在 $D(X\times\Bun_G)$ 中都有函子性的同构
\[H_V(\cF)\simeq V_E\boxtimes\cF,\]
并且这些同构与 $\Rep(\checkG)$ 中的张量积相容，则称对象 $\cF\in D(\Bun_G)$ 是以 $E_{\checkG}$ 为特征值的 Hecke 特征层。特别地，在 $x\in X$ 处逐纤维有
\[H_{V,x}(\cF)\simeq(V_E)_x\otimes\cF.\]
\end{definition}
\srcpage{15}

这是在层论层面替代经典条件：一个自守函数应是所有球面 Hecke 算子的特征向量。
\section{\texorpdfstring{$\GL_n$}{GLₙ} 的微小 Hecke 算子}\label{sec:2.7}\index{微小 Hecke 算子|hyperpage}
对 $G=\GL_n$，有 $\checkG=\GL_n$。令 $\mathrm{Std}$ 为 $\checkG$ 的标准表示。外幂 $\bigwedge^i\mathrm{Std}$，$0\le i\le n$，在几何 Satake 下对应于微小 Schubert 轨道 $\Gr_G^i\subset\Gr_G$。这些是与标准 Hecke 算子有关的闭球面轨道，给出以 $0\le i\le n$ 为指标的 $n+1$ 个标准算子。若保留 $\GL_n$ 的中心次数移位，则可以用中心余权平移这些轨道；模掉这一中心方向后，$\omega_0$ 与 $\omega_n$ 被等同，因而有 $n$ 个类。

令 $\Hecke_n^i$ 为相应的整体 Hecke 对应，$\Gr_{n,X}^i$ 为 $\Gr_{\GL_n,X}$ 中相应的相对 Schubert 分层。一个 $S$-点是向量丛在 $x:S\to X$ 处的修改，其局部相对位置为 $\omega_i$；局部而言，这等价于一个支撑在 $\Gamma_x$ 上、长度为 $i$ 的商。局部模型是通常的 Grassmann 簇 $\Gr(i,n)$，且
\[\Hecke_n^i\simeq\Gr_{n,X}^i\widetilde\times\Bun_n.\]

\begin{transnote}{含糊点拨}
相对位置 $\omega_i=(1,\ldots,1,0,\ldots,0)$ 是 $L^+G$-轨道的标号：在一点取基后，修改映射在圆盘上可写成 $t^{\omega}$ 的矩阵形（Birkhoff 分解），$\omega_i$ 恰对应「余核为 $i$ 维」的那条轨道。$\Hecke_n^i$ 只取这一条，所以一个 $S$-点就是 $\mathcal V'\subset\mathcal V$ 而商支撑在 $\Gamma_x$ 上、长度 $i$；在代数闭点上商是 $k^{\oplus i}$，长度为 $i$ 的商的参数空间与 $i$ 维子空间的参数空间（即 $\Gr(i,n)$）经取对偶相同——这就是「局部模型是通常的 Grassmann 簇」的确切含义。整条轨道给出的参数空间是 $\Gr_{n,X}^i$，于是 $\Hecke_n^i\simeq\Gr_{n,X}^i\widetilde\times\Bun_n$ 逐点地看就是「底丛加上这个参数」。
[来源：nLab "affine Grassmannian"，https://ncatlab.org/nlab/show/affine+Grassmannian（Last revised 2023-08-15；2026-11-17 访问）；并见原书 p.15 对 $\omega_i$ 与 $\Gr(i,n)$ 的说明。]
\end{transnote}
相应的 Hecke 函子为
\[H_n^i:D(\Bun_n)\longrightarrow D(X\times\Bun_n),\]
\[H_n^i(\cF)=(\supp,\overrightarrow h^i)_!\overleftarrow h^{i,*}\cF[\text{移位}],\]
其中移位采用 Satake／反常归一化；对微小轨道，可以取移位 $\dim\Gr(i,n)=i(n-i)$，在 $\ell$-进归一化中配上相应的 Tate 扭转。

\begin{transnote}{易错预警}
两个因子各司其职，不能只取其一。移位 $[d_i]$ 取相对维数 $d_i=i(n-i)=\dim\Gr(i,n)$，使 $\Gr(i,n)$ 上的常值层经平移后落在反常层的 $t$-结构中；Tate 扭转 $(d_i/2)$ 再把 Frobenius 权拉回零（$\mathbb Q_\ell[d_i](d_i/2)$ 是权 $0$ 的）。两者合起来才使下面 (2.1) 的 $H_n^i(\cF)\simeq\bigwedge^iE\boxtimes\cF$ 两侧归一化一致：只移不扭，$\ell$-进情形会多出 $q_x^{d_i/2}$ 的因子；只扭不移，Hecke 函子不再 $t$-正合。
[来源：nLab "perverse sheaf"，https://ncatlab.org/nlab/show/perverse+sheaf（Last revised 2026-10-16；2026-11-17 访问）；nLab "Tate twist"，https://ncatlab.org/nlab/show/Tate+twist（Last revised 2025-08-10；2026-11-17 访问）；并见原书 p.16 式 (2.1) 的归一化。]
\end{transnote}

于是，以秩 $n$ 局部系统 $E$ 为特征值的 Hecke 特征层 $\cF$ 满足
\begin{equation}\label{eq:2.1}
H_n^i(\cF)\simeq\bigwedge^iE\boxtimes\cF,\qquad0\le i\le n.
\end{equation}
其余小节将对由 $E$ 构造的自守层证明这个公式。

现在给出 $\GL_n$ 情形的详细论证。关键是：第一 Hecke 特征性质，加上交换不同点时通常的对称性，就能推出式\eqref{eq:2.1}中所有微小特征性质。这是 Frobenius 特征值的初等对称函数关系的层论形式。
\section{微小 Hecke 对应}\label{sec:2.8}
令 $\Bun_n:=\Bun_{\GL_n}$。对 $0\le i\le n$，置 $\omega_i=(1^i,0^{n-i})$，从而 $t^{\omega_i}=\operatorname{diag}(tI_i,I_{n-i})$，并令
\[\Gr_n^i:=L^+\GL_n\cdot t^{\omega_i}\cdot L^+\GL_n/L^+\GL_n.\]
令 $\Gr_{n,X}^i$ 为 $X$ 上相应的相对 Schubert 分层。相应的整体 Hecke 叠为
\[\Hecke_n^i\simeq\Gr_{n,X}^i\widetilde\times\Bun_n,\]
并具有映射
\origdiagram{p015-diagram.pdf}{0.4303}
一个 $S$-点分类点 $x:S\to X$ 和向量丛的嵌入 $\mathcal V'\hookrightarrow\mathcal V$，使 $\mathcal V/\mathcal V'$ 是支撑于 $\Gamma_x$ 上、相对位置为 $\omega_i$ 的长度 $i$ 挠层。在一个代数闭点上，这一条件是 $\mathcal V/\mathcal V'\simeq k(x)^{\oplus i}$；等价地，$\mathcal V'\subset\mathcal V\subset\mathcal V'(x)$ 且 $\length(\mathcal V/\mathcal V')=i$。
\srcpage{16}

置 $d_i=i(n-i)$。这是 $\Hecke_{n,x}^i\to\Bun_n$ 的相对维数。归一化的 Hecke 函子为
\[H_n^i(\cF):=(\supp,\overrightarrow h)_!\overleftarrow h^*\cF[d_i](d_i/2)\in D(X\times\Bun_n).\]
以 $E$ 为特征值的第 $i$ 个 Hecke 特征性质是
\[H_n^i(\cF)\simeq\bigwedge^iE\boxtimes\cF.\]
\section{轨道记号与闭包序}\label{sec:2.9}\index{Schubert 轨道|hyperpage}\index{优势序|hyperpage}
整体 Hecke 叠通过用 $X\times\Bun_G$ 上的无限级 $L^+G$ 挠子扭曲仿射 Grassmann 簇而得到：
\[\Hecke_{G,X}\simeq\Gr_{G,X}\widetilde\times\Bun_G.\]
因此，每个 $L^+G$ 轨道 $\Gr_G^\lambda\subset\Gr_G$ 都给出相应的相对分层 $\Gr_{G,X}^\lambda\subset\Gr_{G,X}$，以及
\[\Hecke_G^\lambda\simeq\Gr_{G,X}^\lambda\widetilde\times\Bun_G.\]
对 $G=\GL_n$，优势余权为元组 $\lambda=(\lambda_1\ge\cdots\ge\lambda_n)$，且
\[\Gr_n^\lambda:=L^+\GL_n\cdot t^\lambda\cdot L^+\GL_n/L^+\GL_n,\qquad t^\lambda=\operatorname{diag}(t^{\lambda_1},\ldots,t^{\lambda_n}).\]
在一个固定的连通分支内，闭包序就是优势序：$\Gr_n^\lambda\subset\overline{\Gr_n^\mu}$ 当且仅当对所有 $1\le r<n$ 都有 $\sum_{j=1}^r\lambda_j\le\sum_{j=1}^r\mu_j$。微小余权 $\omega_i=(1^i,0^{n-i})$ 在几何 Satake 下对应于 $\bigwedge^i\mathrm{Std}$。
\section{经典自守归一化}\label{sec:2.10}
令
\[G_n^i=L^+\GL_n\cdot\operatorname{diag}(\underbrace{t,\ldots,t}_{i},1,\ldots,1)\cdot L^+\GL_n/L^+\GL_n.\]
对球面函数 $f\in C_c^\infty(L^+\GL_n\backslash L\GL_n/L^+\GL_n)$，归一化的局部 Hecke 算子为
\[T_x^i(f)(y)=q_x^{-d_i/2}\int_{G_n^i}f(yg)\,dg.\]
若局部系统 $E$ 对应 Weil 表示 $\rho:W_F\to\GL_n(\mathbb Q_\ell)$，则预期特征值为 $\operatorname{tr}(\operatorname{Frob}_x;\bigwedge^i\rho)$。等价地，
\[\det(1-\rho(\operatorname{Frob}_x)s)=\sum_{i=0}^n(-1)^i\operatorname{tr}(\operatorname{Frob}_x;\bigwedge^i\rho)s^i.\]
这是 $H_n^i(\cF)\simeq\bigwedge^iE\boxtimes\cF$ 在函数论中的体现。
\section{归约到第一 Hecke 算子}\label{sec:2.11}
假设 $\cF\in\operatorname{Perv}(\Bun_n)$ 配备同构
\[\alpha:H_n^1(\cF)\xrightarrow{\sim}E\boxtimes\cF.\]
对两个不同的点 $x\ne x'$，两次初等修改的叠有 $S_2$ 作用：考虑 $\mathcal V''\subset\mathcal V'\subset\mathcal V$，满足 $\length(\mathcal V/\mathcal V')=\length(\mathcal V'/\mathcal V'')=1$、$\supp(\mathcal V/\mathcal V')=x$、$\supp(\mathcal V'/\mathcal V'')=x'$。交换 $x$ 与 $x'$ 就将这两次初等修改的次序颠倒。特征层所要求的相容性是同构
\[H_n^1H_n^1(\cF)|_{X^2-\Delta}\xrightarrow{\sim}E\boxtimes E\boxtimes\cF|_{X^2-\Delta}\]
为 $S_2$-等变的。迭代后，在 $X^i-\Delta$ 上得到 $S_i$-等变同构
\[(H_n^1)^i(\cF)|_{X^i-\Delta}\simeq E^{\boxtimes i}\boxtimes\cF|_{X^i-\Delta}.\]
\srcpage{17}
\begin{proposition}{}{2.11.1}
在上述对称性假设下，对每个 $0\le i\le n$，都有
\[H_n^i(\cF)\simeq\bigwedge^iE\boxtimes\cF.\]
\end{proposition}
\section{几何 Satake 的输入}\label{sec:2.12}
回顾
\[\operatorname{Perv}_{L^+G}(\Gr_G)\simeq\operatorname{Perv}(L^+G\backslash LG/L^+G)\simeq\Rep(\checkG),\]
沿这一等价，对所有 $\lambda$ 都有
\[\operatorname{IC}_{\overline{\Gr_G^\lambda}}\longleftrightarrow V^\lambda.\]
对 $G=\GL_n$，这给出
\[\mathrm{Std}\longleftrightarrow(1,0,\ldots,0)=(1,0^{n-1}),\qquad\bigwedge^i\mathrm{Std}\longleftrightarrow(1^i,0^{n-i}).\]
因此，$\operatorname{IC}_{\Gr_n^1}^{\star i}$ 对应于 $\mathrm{Std}^{\otimes i}$，自然 $S_i$ 作用的交错投影算子切出 $\bigwedge^i\mathrm{Std}$。整体证明通过长度 $i$ 修改叠的 Springer 型解消来实现这一投影算子。
\section{长度 \texorpdfstring{$i$}{i} 修改的紧化叠}\label{sec:2.13}\index{Springer 解消|hyperpage}\index{Springer 层|hyperpage}
令
\[\operatorname{Mod}_n^{-i}:=\{(\mathcal V'\hookrightarrow\mathcal V)\mid\mathcal V/\mathcal V'\in\Tor^i\}.\]
它有对应
\origdiagram{p017-diagram.pdf}{0.4258}
其中 $\supp(\mathcal V/\mathcal V')$ 是相应的有效除子。令 $\Delta_i:X\to X^{(i)}$ 为 $x\mapsto i\cdot x$，并置
\[H_n^{i,+}:=X\times_{X^{(i)}}\operatorname{Mod}_n^{-i}.\]
满足 $\mathcal V/\mathcal V'\simeq k(x)^{\oplus i}$ 的层正是 $\Hecke_n^i$。因此，$H_n^{i,+}$ 通过允许支撑在 $x$ 处的非半单长度 $i$ 商，将微小对应紧化。

对 $I=\{1,\ldots,i\}$，Beilinson–Drinfeld Grassmann 簇给出
\[\Gr_{G,X}^I\widetilde\times\Bun_G=\Hecke_{G,X}^I.\]
对 $G=\GL_n$，存在从正的整体 Schubert 簇到紧化修改叠的真态射：
\[\Gr_{G,X}^{+,\le(1,\ldots,1)}\widetilde\times\Bun_n=\Hecke_{n,X}^{+,\le(1,\ldots,1)}\longrightarrow\operatorname{Mod}_n^{-i}.\]
在 $X^i-\Delta$ 上，它因子化为
\[\Gr_{G,X}^{+,\le(1,\ldots,1)}|_{X^i-\Delta}\simeq(\Gr_{G,x_1}^1\times\cdots\times\Gr_{G,x_i}^1)|_{X^i-\Delta}.\]

现在定义有序解消，也称 Springer 解消：
\[\widetilde{\operatorname{Mod}}_n^{-i}:=\{\mathcal V_0\hookrightarrow\mathcal V_1\hookrightarrow\cdots\hookrightarrow\mathcal V_i\mid\length(\mathcal V_j/\mathcal V_{j-1})=1\text{，对所有 }j\}.\]
存在真态射
\[\varphi:\widetilde{\operatorname{Mod}}_n^{-i}\longrightarrow\operatorname{Mod}_n^{-i},\qquad(\mathcal V_0\subset\cdots\subset\mathcal V_i)\longmapsto(\mathcal V_0\subset\mathcal V_i),\]
\srcpage{18}
以及有序支撑映射 $r:\widetilde{\operatorname{Mod}}_n^{-i}\to X^i$。它们组成交换图
\origdiagram{p018-square.pdf}{0.2755}
在互异点轨迹上，$\varphi$ 是有序支撑覆盖，并带有通过置换长度一商的次序给出的自然 $S_i$ 作用。
\begin{claim}{}{2.13.1}
Springer 层
\[\operatorname{Spr}_i:=\varphi_!\operatorname{IC}_{\widetilde{\operatorname{Mod}}_n^{-i}}\]
满足以下性质：
\begin{enumerate}
\item $\operatorname{Spr}_i$ 是反常层。
\item $\operatorname{Spr}_i$ 是从互异点轨迹作出的中间延拓；在该轨迹上，它是与有序支撑覆盖及 $S_i$ 的正则表示相联系的局部系统。
\item 交错直和项是微小核：
\[\Hom_{S_i}(\mathrm{sgn},\operatorname{Spr}_i)|_{\Hecke_n^i}\simeq\mathbb Q_\ell|_{\Hecke_n^i}[d_i](d_i/2),\qquad d_i=i(n-i).\]
\end{enumerate}
\end{claim}

这与通常 Springer 层的机制相同：解消为初等商排序，符号表示则提取外幂直和项。
\section{投影公式计算}\label{sec:2.14}
\enlargethispage{\baselineskip}
\begin{proof}
通向 $\Bun_n$ 的解消奇点映射为
\origdiagram{p018-triangle.pdf}{0.4258}
支撑映射组成如下支撑方块：
\origdiagram{p018-square.pdf}{0.2755}
由投影公式，
\[(\supp,\overrightarrow h)_!(\overleftarrow h^*\cF\otimes\operatorname{Spr}_i)\simeq(\mathrm{sym}\times\mathrm{id}_{\Bun_n})_!((H_n^1)^i(\cF)).\]
这里，归一化的移位与 Tate 扭转已经包含在 $\operatorname{Spr}_i$ 的 IC 归一化以及函子 $H_n^1$ 中。在 $X^i-\Delta$ 上，右边是迭代第一 Hecke 函子的对称化。利用 $S_i$-等变的特征同构以及中间延拓性质，得到
\[(\mathrm{sym}\times\mathrm{id}_{\Bun_n})_!((H_n^1)^i(\cF))\simeq\mathrm{sym}_!(E^{\boxtimes i})\boxtimes\cF.\]
取符号表示的等型直和项，再沿 $\Delta_i:X\to X^{(i)}$ 限制，便得到
\[\Delta_i^*\Hom_{S_i}(\mathrm{sgn},\mathrm{sym}_!E^{\boxtimes i})\simeq\bigwedge^iE.\]
\srcpage{19}
将这个恒等式与核的符号 Springer 描述结合，就得到
\[H_n^i(\cF)\simeq\bigwedge^iE\boxtimes\cF.\]
这就证明了命题\ref{pro:2.11.1}。
\end{proof}
